\section{D-frame and Q-enriched refinements}

\begin{proposition}[D-frames are the algebraic dual carrier]\label{thm:dframe-carrier}\lean{Infomaterial.DFrame.DFrame}\leanok\uses{thm:bitopological-site-grothendieck}
A d-frame carries two frames together with consistency and totality relations.
It is the algebraic dual surface for the bitopological settling site.
\end{proposition}

\begin{theorem}[The bitopological d-frame has disjointness and covering relations]\label{thm:bitopological-dframe-relations}\lean{Infomaterial.DFrame.dframe_of_bitopological_relations}\leanok\uses{thm:dframe-carrier, thm:four-priestley-dual}
For a bitopological point set, consistency is disjointness of the positive and
negative opens, and totality is covering.  These are the spatial relations the
settling sheaf refines.
\end{theorem}

\begin{theorem}[The settling sheaf is d-frame-valued]\label{thm:dframe-valued-sheaf}\lean{Infomaterial.DFrame.dframe_valued_sheaf}\leanok\uses{thm:bitopological-dframe-relations, thm:settling-associated-sheaf}
The d-frame-valued settling sheaf packages the truth side and the
falsehood/inconsistency side at once.  Restriction along a clopen inclusion is
a d-frame morphism, preserving meet, top, monotonicity, and consistency.
\end{theorem}

\begin{proposition}[The FOUR is a charge quantale]\label{thm:charge-quantale}\lean{Infomaterial.QEnriched.chargeQuantale}\leanok\uses{thm:four-bilattice}
The charge lattice is a commutative unital quantale: multiplication is
knowledge meet, the unit is knowledge top, and suprema are knowledge joins.
\end{proposition}

\begin{theorem}[Q-categories form over the charge quantale]\label{thm:qcat-is-quantale-category}\lean{Infomaterial.QEnriched.qcat_is_quantale_category}\leanok\uses{thm:charge-quantale}
The one-object and discrete charge-enriched categories satisfy the Q-category
laws, and Q-functors have identity and composition.
\end{theorem}

\begin{theorem}[Q-fuzzy equivalence composes by multiplying grades]\label{thm:q-fuzzy-equivalence-composes}\lean{Infomaterial.QEnriched.q_fuzzy_equivalence_composes}\leanok\uses{thm:qcat-is-quantale-category}
Q-fuzzy equivalences compose with multiplicative grade, and grade \(1\)
recovers the strict equivalence bound.  This is the enriched comparison
surface for later q-fuzzy cohomology claims.
\end{theorem}

\begin{theorem}[Degree-zero Q-cohomology is Laplacian fixed points]\label{thm:h0-q-laplacian-fixed-points}\lean{Infomaterial.QEnriched.h0_laplacian_fixed_points}\leanok\uses{thm:q-fuzzy-equivalence-composes, thm:settling-associated-sheaf}
In the charge-quantale degree-zero case, the Q-enriched Laplacian is knowledge
meet against a self-hom weight, and \(H^0_Q\) is its fixed-point carrier.
\end{theorem}
